
\newpage
% \color{white}
% \pagecolor{black}

\ifdefined\useChapters
\chapter{Sozinho}
\else
\chapter{}
\fi

Na manhã seguinte à nossa volta pra minha casa, logo cedo, descemos pra tomar café e eu finalmente ia apresentar o Ítalo pra minha irmã. Ela precisava ver que eu podia ter amigos de novo. Desde que despertei do coma, ela estava muito preocupada porque eu sempre andava sozinho e por isso já tinha me falado mais de dez vezes que eu estava estranho e ia me levar a um psicólogo.

Enquanto descíamos as escadas, o cheiro de ovos fritos já inundava toda a casa. Minha irmã era uma ótima cozinheira e sabia muito bem como me mimar. Quando chegamos na cozinha, a mesa já estava arrumada. Em cima dela tinham três pratos, uma jarra de suco de laranja, algumas torradas e ovos fritos.

— Bom dia, meu irmão. Achei que nunca mais fosse te ver de novo.

— Oi, Lúcia, que saudade.

Ela me olhou furiosa e me deu um soco no braço.

— Eu também, demais — falou ela enquanto me abraçava tão forte que até senti uma pequena falta de ar.

— Onde você estava todos esses dias? A diretora me ligou várias vezes dizendo que você não estava indo à escola.

— Não vai me apresentar seu amigo? — perguntou ela ainda abraçada comigo, olhando pra ele por sobre os meus ombros.

— Ah, esse é o Ítalo, ele é meu amigo da escola.

— Amigo da escola? E ele mora onde? Eu nunca o vi aqui por perto — disse desconfiada enquanto me soltava.

— Eu sou novo na escola, morava no interior e meus pais se mudaram pra cá faz alguns meses — respondeu Ítalo, sustentando a história para que ela não desconfiasse.

— Bom, seja bem-vindo. Meu irmão é meio avoado, sempre tem umas ideias meio doidas, mas no fundo é gente boa.

Sentamos à mesa e ela começou a fazer perguntas pra tentar conhecer meu novo amigo. Ítalo é muito esperto; eu quase me convenci de que nos conhecíamos bem realmente.

— Muito obrigado pelo café, mas a gente está super atrasado. Outra hora te explico um pouco melhor onde eu estive todos esses dias.

— Atrasado pra quê? O que é mais importante do que se explicar pra mim?

Enquanto ela falava, olhei pra Ítalo como quem diz: vamos antes que ela faça mais perguntas. Ele, entendendo completamente, pegou umas torradas da mesa, enfiou na boca e me puxou pela mochila, saindo pela cozinha.

— A gente não tem muito tempo, cara — falou ele, e assim que pegamos uma distância da minha casa, me segurou com força e começou a volitar muito rápido pra cima.

— O que é isso, cara? As pessoas vão nos notar assim.

— Notar? Olha pra baixo, todo mundo está tão preocupado com suas vidas que nunca olham pra cima. Eles têm certeza de que nunca vai ter nada diferente pra ver.

— Aliás, você precisa aprender logo a dominar suas habilidades. Já passou da hora de eu não ter que ficar te carregando pra cima e pra baixo. Se continuar assim, não vamos ter nenhuma chance contra eles. Como eu já te disse várias vezes, além de você, ninguém mais está brincando de superpoderes por aqui.

Nos próximos minutos não falamos nada e ele só me carregou por sobre a cidade. Eu não tinha muito o que dizer, ele estava com toda a razão. Não posso mais me dar ao luxo de ficar com medo de usar as habilidades ou usá-las só para satisfazer minhas vontades.

— Chegamos — falou Ítalo enquanto me soltava no chão.

— Onde você me trouxe?

— Te trouxe pra montanha mais alta da cidade e não vai sair daqui enquanto você não me mostrar do que é capaz — disse Ítalo enquanto volitava, lá do alto me acenando como se estivesse batendo continência.

— Volta aqui, cara, você não pode me deixar nesse lugar.

— Tanto posso que já deixei, te vejo lá embaixo.

Eu não acredito que ele fez isso, me trouxe pra uma montanha alta pra caramba e, se não bastasse isso, ela é tão íngreme que eu nunca vou conseguir descer daqui. Eu sabia que ele não era meu amigo, eu sou mesmo muito idiota pra confiar, sempre tive esse problema, sempre acredito nas pessoas.

— Como eu sou... otário!

Não sei por que estou gritando, daqui de cima aposto que ninguém consegue me ouvir e o pior de tudo é que já está começando a escurecer, preciso buscar um abrigo.

Estou preso aqui em cima, sozinho, com fome e frio.

Eu não posso me render a isso, preciso me acalmar pra pensar melhor. Sento então com as pernas cruzadas e fecho meus olhos.

O que eu posso fazer?

Não sei escalar e, mesmo que soubesse, com certeza morreria caindo daqui de cima.

Se eu ficar aqui, vou morrer de fome e frio.

Eu preciso volitar e sair daqui, mas eu não consigo me lembrar como faço isso. Das últimas vezes consegui porque não tinha outra alternativa, e é o caso agora.

Na posição em que eu estava, sem perceber caí no sono e despertei algumas horas depois com o sol esquentando minha cabeça, um vento muito forte na cara e o barulho da minha barriga roncando.

— Maldito, cadê você?

Quem aquele cara pensa que é pra me prender aqui? Assim que eu sair daqui, vou pegar ele.

Por que eu não consigo usar minhas habilidades? Que droga.

Preciso encontrar alguma coisa pra comer ou beber, senão vou morrer aqui em cima, mas olhando pro lado, não tem nada por aqui. Ele me trouxe para o topo de uma maldita montanha que deve ter uns cinquenta metros quadrados, com muita grama, pedras e uma árvore que nem pra ser frutífera serve.

De novo me sento com as pernas cruzadas e me concentro em entender o que está acontecendo, por que não consigo usar nenhuma habilidade? De novo, de tanto meditar, pego no sono.

De novo acordo aqui em cima, mas dessa vez consigo, depois de um tempo, perceber como aqui é bonito. É impressionante como o nascer do sol daqui é lindo. Aos poucos começo a sentir um pouco mais o vento batendo no meu rosto e sinto meus pelos do corpo inteiro levantarem. Sinto aos poucos um pouco mais o meu corpo, me percebo um pouco mais e lembro que todas as vezes em que consegui utilizar alguma habilidade, eu estava em um momento de desespero, meus sentidos estavam muito mais aguçados. Eu preciso sentir no meu corpo e acreditar que sou capaz. Nesses últimos dias aqui em cima, só tinha raiva tomando meus pensamentos e mudando o meu foco pra outra pessoa. Preciso sentir e acreditar.

Me levanto, ergo a minha cabeça e chego bem próximo da borda do desfiladeiro. Se eu ficar aqui mais um dia, vou morrer mesmo, então preciso dar um salto de fé.

Abro meus braços e começo a sentir um enorme frio na espinha. Eu vou mesmo fazer isso, vou mesmo me jogar?

Me encho de coragem, respiro bem fundo e me jogo.

Por um tempo, caio da mesma maneira que caí quando me soltei do Ítalo enquanto era carregado.

Caio mais um pouco, e aos poucos, o vento que sinto começa a diminuir, minha velocidade começa a diminuir e eu começo a volitar como já havia feito, mas dessa vez é diferente, porque eu sei o que estou fazendo, não é algo involuntário.

Quando eu tomo total controle e passo a volitar tão rápido quanto Ítalo, indo em direção à cidade, sou abordado.

— Finalmente você conseguiu, chegou a passar pela minha cabeça que você morreria lá em cima — Ítalo me fala enquanto volita próximo a mim.

— Saí de lá, mas não foi graças a você.

— Deixa de ser ingrato, eu fiz isso de propósito, foi a única maneira de você se destravar e entender de uma vez por todas a natureza das habilidades.

Ele tinha me deixado preso lá em cima alguns dias e eu não podia confiar nele assim de novo, do nada. Me lembro de como me senti enquanto lutava com Larissa e das minhas mãos começam a brotar chamas como naquele dia.

— É disso que estou falando — grita Ítalo entusiasmado.

— Você está maluco? Por que você está feliz? Não acha que eu vou te atacar?

— Porque independentemente de você acreditar em mim, que eu fiz o que fiz pra te ajudar, você agora parece mais pronto pra enfrentar o que precisamos.

— Agora pare de sentimentalismo e me siga, a gente precisa pensar em como vamos combater Crispim e os outros — me fala Ítalo como quem sabe que eu vou segui-lo, enquanto me dá as costas e dispara a volitar em direção à cidade.

Ele está certo, não temos tempo pra lutar entre nós, mas agora, caso ele me traia de novo, eu sei muito bem como me defender.

